%
% Este archivo es un ejemplo que probablemente no se necesite.
% Puede borrarse sin ningún efecto en la plantilla.
%
\chapter{Definición del problema}
\label{defProblema}  % para poder enlazar a este capítulo

% Se puede subdividir todo en los ficheros que queramos. Lo normal es por capítulos. No es recomendable subdividir si un capítulo va a ser pequeño pero dejo este así como ejemplo por si os puede interesar para algo. NO lo aconsejo para empezar.

% El uso de input (aquí) o include (que se ha usado en el principal) se diferencia en que con include se compila cada apartado independientemente (puede acelerar la ejecución pero genera más archivos). La división con input es sólo a efectos de organización nuestra para editar (antes de compilar es como si se pegase el fichero del input tal cual en lugar del input).
\input{021_Definicion-Problema-Real}
% \input{022_CapSegundoParteDos}
% \input{023_CapSegundoParteTres}
% \input{024_CapSegundoParteCuatro}

